\subsubsection{第1段階ステップB}
\label{sec:model_households_indices_aggregation_stage1b}

次に第1段階（期内費用最小化）のうちステップBを考える。

\paragraph{自国家計の第1段階ステップB}
\label{sec:model_households_indices_aggregation_stage1b_home}

\subparagraph{総消費指数の定義}
\label{sec:model_households_indices_aggregation_stage1b_home_c_h_W}
総消費指数\(c_t^{h\to W}\)は自国財消費指数\(c_t^{h\to H}\)と外国財消費指数\(c_t^{h\to F}\)から合成される。
これはCES型指数において代替の弾力性パラメータ\(\theta\)が1である特別な場合（コブ＝ダグラス型）に対応する。
総消費指数\(c_t^{h\to W}\)は以下のように定義される。
\begin{equation}
\label{form:model_households_indices_aggregation_stage1b_home_c_h_W_def}
    c_t^{h\to W}\equiv\frac{(c_t^{h\to H})^{\alpha^H}(c_t^{h\to F})^{1-\alpha^H}}{(\alpha^H)^{\alpha^H}(1-\alpha^H)^{1-\alpha^H}}
\end{equation}
この定義はパラメータ\(\alpha^H\)を用いた加重幾何平均の形をとる。
この形式がもつ一つの重要な含意は単位の整合性である。
仮に消費指数の単位を「単位」とすると
\begin{equation}
\label{form:model_households_indices_aggregation_stage1b_home_c_h_W_unit_id}
    (\text{単位})^{\alpha^H}\times(\text{単位})^{1-\alpha^H}=\text{単位}^{\alpha^H+1-\alpha^H}=\text{単位}^1
\end{equation}
となり総消費指数\(c_t^{h\to W}\)はその構成要素と同じ次元を持つ。
これにより\(c_t^{h\to W}\)は家計の総体的な消費水準を示す指標として成立する。
\(1/\left((\alpha^H)^{\alpha^H}(1-\alpha^H)^{1-\alpha^H}\right)\)という定数項は
総消費指数\(c_t^{h\to W}\)を調整するための正規化項である。
モデルの式中では消費者物価指数\(p_t^{H\to W}\)が多用されるため、
その定義式が簡潔になるよう、あらかじめ\(c_t^{h\to W}\)の定義にこの項を入れておくのが一般的である。
この正規化により後続する消費者物価指数\(p_t^{H\to W}\)の定義式から
定数項\((\alpha^H)^{\alpha^H}(1-\alpha^H)^{1-\alpha^H}\)が消え、
\(p_t^{H\to W}=(p_t^H)^{\alpha^H}(p_t^F)^{1-\alpha^H}\)という簡潔な形で表現される。

\subparagraph{期内費用最小化問題とラグランジュ乗数}
次に家計\(h\)は所与の総消費指数\(c_t^{h\to W}\)を達成するために
自国財消費指数\(c_t^{h\to H}\)と外国財消費指数\(c_t^{h\to F}\)の
組み合わせ費用\(p_t^Hc_t^{h\to H}+p_t^Fc_t^{h\to F}\)を最小化する問題を考える。
この問題のラグランジアンは以下のように書ける。
\begin{equation}
\label{form:model_households_indices_aggregation_stage1b_home_lagrangian_def}
    \mathcal{L}^{h\to W}=p_t^Hc_t^{h\to H}+p_t^Fc_t^{h\to F}+\eta_t^h\left(c_t^{h\to W}-\frac{(c_t^{h\to H})^{\alpha^H}(c_t^{h\to F})^{1-\alpha^H}}{(\alpha^H)^{\alpha^H}(1-\alpha^H)^{1-\alpha^H}}\right)
\end{equation}
ここでラグランジュ乗数\(\eta_t^h\)は
家計\(h\)にとっての総消費指数\(c_t^{h\to W}\)の限界価格を表す。
この費用最小化問題の1階条件から各消費指数への需要関数は\(\eta_t^h\)を用いて以下のように導かれる。
\begin{align}
    \label{form:model_households_indices_aggregation_stage1b_home_p_H_c_h_H_eq}
        p_t^Hc_t^{h\to H}&=\alpha^H\eta_t^hc_t^{h\to W} \\
    \label{form:model_households_indices_aggregation_stage1b_home_p_F_c_h_F_eq}
        p_t^Fc_t^{h\to F}&=(1-\alpha^H)\eta_t^hc_t^{h\to W}
\end{align}
この結果を名目総消費\(p_t^Hc_t^{h\to H}+p_t^Fc_t^{h\to F}\)に代入すると
\begin{equation}
\label{form:model_households_indices_aggregation_stage1b_home_total_nominal_consumption_eq}
    p_t^Hc_t^{h\to H}+p_t^Fc_t^{h\to F}=\alpha^H\eta_t^hc_t^{h\to W}+(1-\alpha^H)\eta_t^hc_t^{h\to W}=\eta_t^hc_t^{h\to W}
\end{equation}
という関係が得られる。
これは第1段階と同様に\(\eta_t^h\)が総消費指数の限界価格であることから予想された結果である。

\subparagraph{消費者物価指数(CPI)の定義}
この最適化問題を解く過程で限界価格\(\eta_t^h\)は以下の形で決定される
（詳細は付録\ref{chap:appendix_cost_minimization}参照）。
\begin{equation}
\label{form:model_households_indices_aggregation_stage1b_home_eta_h_eq}
    \eta_t^h=(p_t^H)^{\alpha^H}(p_t^F)^{1-\alpha^H}
\end{equation}
この式からわかるように\(\eta_t^h\)は各国の生産者物価指数\(p_t^H,p_t^F\)のみで決定され
家計\(h\)自身の選択には依存しない。
したがって\(\eta_t^h\)はすべての家計で共通の値をとる。
この共通の限界価格を\(p_t^{H\to W}\)と定義し本稿では消費者物価指数(CPI)と呼ぶ。
\begin{equation}
\label{form:model_households_indices_aggregation_stage1b_home_p_H_W_def}
    p_t^{H\to W}\equiv\eta_t^h=(p_t^H)^{\alpha^H}(p_t^F)^{1-\alpha^H}
\end{equation}
これとは別にマクロ分析のため経済規模の影響(\(N,M\))を取り除いた
「正規化された」消費者物価指数\(\bar{p}_t^{H\to W}\)を次のように定義する。
\begin{equation}
\label{form:model_households_indices_aggregation_stage1b_home_p_H_W_bar_def}
    \bar{p}_t^{H\to W}\equiv(\bar{p}_t^H)^{\alpha^H}(\bar{p}_t^F)^{1-\alpha^H}
\end{equation}
この\(\bar{p}_t^{H\to W}\)は個々の企業の価格を直接集計した統計上の消費者物価指数に相当する。
中央銀行が消費者物価指数インフレ目標を採用する場合、
目標とするインフレ率はこの正規化された消費者物価指数\(\bar{p}_t^{H\to W}\)の変化率として算出される。

\subparagraph{自国家計の需要関数と支出シェア}
この関係\(\eta_t^h=p_t^{H\to W}\)を式
\eqref{form:model_households_indices_aggregation_stage1b_home_p_H_c_h_H_eq}と
\eqref{form:model_households_indices_aggregation_stage1b_home_p_F_c_h_F_eq}に代入し、
\(c_t^{h\to H}\)と\(c_t^{h\to F}\)について解くことで需要関数の最終形が得られる。
\begin{align}
    \label{form:model_households_indices_aggregation_stage1b_home_c_h_H_eq}
        c_t^{h\to H}&=\alpha^H\frac{p_t^{H\to W}}{p_t^H}c_t^{h\to W} \\
    \label{form:model_households_indices_aggregation_stage1b_home_c_h_F_eq}
        c_t^{h\to F}&=(1-\alpha^H)\frac{p_t^{H\to W}}{p_t^F}c_t^{h\to W}
\end{align}
ここで定義上は集計の際の重みとして導入されたパラメータ\(\alpha^H\)が
家計の最適化の結果として支出シェアという経済学的な意味を持つことを示す。
式\eqref{form:model_households_indices_aggregation_stage1b_home_p_H_c_h_H_eq}は
自国財消費指数への名目支出額\(p_t^Hc_t^{h\to H}\)が総消費指数への名目総支出額
\(\eta_t^hc_t^{h\to W}=p_t^{H\to W}c_t^{h\to W}\)の\(\alpha^H\)倍に等しくなることを
直接的に示している。
同様に式\eqref{form:model_households_indices_aggregation_stage1b_home_p_F_c_h_F_eq}は
外国財消費指数への名目支出額\(p_t^Fc_t^{h\to F}\)が
名目総支出額の\((1-\alpha^H)\)倍になることを示している。
したがって家計が最適な消費の組み合わせを選択した結果、
定義における重みパラメータ\(\alpha^H\)と\(1-\alpha^H\)は
必然的に総支出に占める自国財消費指数と外国財消費指数それぞれへの支出シェアと一致する。
また名目総消費と指数の間には
\(p_t^Hc_t^{h\to H}+p_t^Fc_t^{h\to F}=p_t^{H\to W}c_t^{h\to W}\)という関係が成り立つ。

\paragraph{外国家計の第1段階ステップB}
\label{sec:model_households_indices_aggregation_stage1b_foreign}

\subparagraph{外国家計の消費者物価指数(CPI)}
同様に外国家計\(f\)は所与の総消費指数\(c_t^{f\to W}\)を達成するために
外国財消費指数\(c_t^{f\to F}\)と自国消費指数\(c_t^{f\to H}\)の組み合わせ費用
\(p_t^{F*}c_t^{f\to F}+p_t^{H*}c_t^{f\to H}\)を最小化する。
この結果、外国の消費者物価指数(CPI,外国通貨建て)\(p_t^{F\to W*}\)が
ラグランジュ乗数\(\eta_t^f\)と等しいものとして以下のように定義される。
\begin{equation}
\label{form:model_households_indices_aggregation_stage1b_foreign_p_F_W_star_def}
    p_t^{F\to W*}\equiv\eta_t^f=(p_t^{F*})^{\alpha^F}(p_t^{H*})^{1-\alpha^F}
\end{equation}
これとは別にマクロ分析のため経済規模の影響(\(N,M\))を取り除いた
「正規化された」外国の消費者物価指数\(\bar{p}_t^{F\to W*}\)を次のように定義する。
\begin{equation}
\label{form:model_households_indices_aggregation_stage1b_foreign_p_F_W_star_bar_def}
    \bar{p}_t^{F\to W*}\equiv(\bar{p}_t^{F*})^{\alpha^F}(\bar{p}_t^{H*})^{1-\alpha^F}
\end{equation}
この\(\bar{p}_t^{F\to W*}\)は個々の企業の価格を直接集計した統計上の消費者物価指数に相当する。
中央銀行が消費者物価指数インフレ目標を採用する場合、
目標とするインフレ率はこの正規化された消費者物価指数\(\bar{p}_t^{F\to W*}\)の変化率として算出される。

\subparagraph{外国家計の需要関数}
また外国家計の各消費指数への需要関数は以下の通りである。
\begin{align}
    \label{form:model_households_indices_aggregation_stage1b_foreign_c_f_F_eq}
        c_t^{f\to F}&=\alpha^F\frac{p_t^{F\to W*}}{p_t^{F*}}c_t^{f\to W} \\
    \label{form:model_households_indices_aggregation_stage1b_foreign_c_f_H_eq}
        c_t^{f\to H}&=(1-\alpha^F) \frac{p_t^{F\to W*}}{p_t^{H*}}c_t^{f\to W}
\end{align}
名目総消費と指数の間には\(p_t^{F*}c_t^{f\to F}+p_t^{H*}c_t^{f\to H}=p_t^{F\to W*}c_t^{f\to W}\)という関係が成り立つ。